% Created 2026-03-05 Thu 16:16
% Intended LaTeX compiler: pdflatex
\documentclass[11pt]{article}
\usepackage[utf8]{inputenc}
\usepackage[T1]{fontenc}
\usepackage{graphicx}
\usepackage{longtable}
\usepackage{wrapfig}
\usepackage{rotating}
\usepackage[normalem]{ulem}
\usepackage{amsmath}
\usepackage{amssymb}
\usepackage{capt-of}
\usepackage{hyperref}
\author{The Prologos Project - from Zachary Larson}
\date{2026-03-05}
\title{Protocols as Types: Composable Session Types in Prologos}
\hypersetup{
 pdfauthor={The Prologos Project - from Zachary Larson},
 pdftitle={Protocols as Types: Composable Session Types in Prologos},
 pdfkeywords={},
 pdfsubject={},
 pdfcreator={Emacs 30.2 (Org mode 9.7.39)}, 
 pdflang={English}}
\begin{document}

\maketitle
\tableofcontents

\section{The Insight}
\label{sec:org1f97f50}

Session types are types. Continuations are type positions. Named session types
can appear in continuation positions. Therefore: \textbf{protocols compose through
the type system with no additional mechanism.}

A session type describes a communication protocol. After each operation (send,
receive, select, offer), there is a continuation --- a type describing "what
comes next." When that continuation is a named session type, the protocol
\textbf{transitions into a different protocol}. The type checker verifies the
transition. The channel remains the same; only the protocol changes.

This is not syntactic sugar, not a macro expansion, not a runtime dispatch. It
is the natural consequence of session types being first-class types in a
dependently-typed language. Composition falls out of the type system for free.
\section{Composition Patterns}
\label{sec:orgcd3a82a}

Three composition patterns emerge from this insight, none requiring new
mechanisms beyond what session types already provide.
\subsection{Pattern 1: Phase Transition (Sequential Composition)}
\label{sec:org1261dbb}

A session type's continuation names the next phase of the protocol:

\begin{verbatim}
session Handshake
  ! Version . ? Version . Authenticated  ;; after handshake → authenticate

session Authenticated
  ! Credentials
  ? [Result Token AuthError]
  +>
    | :ok   -> Ready             ;; success → operational phase
    | :fail -> end               ;; failure → done

session Ready
  +>
    | :query -> ! String . ? [Result Rows DbError] . Ready
    | :close -> end
\end{verbatim}

The protocol flows: \texttt{Handshake → Authenticated → Ready → (loop) → end}. Each
phase is a named type. Each transition is verified by the type checker. The
\textbf{author} of each session type decides what comes next --- \texttt{Handshake} declares
"after me comes \texttt{Authenticated}," not the caller.
\subsection{Pattern 2: Mode Branching}
\label{sec:orgc116149}

Choice operators (\texttt{+>} / \texttt{\&>}) can branch into different sub-protocols:

\begin{verbatim}
session DatabaseSession
  +>
    | :query  -> QueryMode
    | :admin  -> AdminMode
    | :close  -> end

session QueryMode
  ! String . ! [List Param] .
  ? [Result [List Row] DbError] .
  DatabaseSession                    ;; return to main protocol

session AdminMode
  +>
    | :vacuum -> ? [Result Unit DbError] . DatabaseSession
    | :backup -> ! Path . ? [Result Unit DbError] . DatabaseSession
\end{verbatim}

This is \textbf{mutual recursion between protocols}. \texttt{DatabaseSession} branches to
\texttt{QueryMode}, which does its work, then transitions back to \texttt{DatabaseSession}.
The protocol graph can be arbitrarily complex --- any directed graph of named
session types is expressible.
\subsection{Pattern 3: Inline Prefix with Named Continuation}
\label{sec:org4b2223c}

Operations are written inline, with a named type taking over as the final
continuation:

\begin{verbatim}
session SecureFileSession
  ! Version . ? Version .                      ;; inline handshake operations
  ! Credentials . ? [Result Token AuthError] . ;; inline auth operations
  FileOps                                      ;; named type takes over
\end{verbatim}

This is the most direct form of "do these things, then become that protocol."
The inline prefix is protocol-specific; the named continuation is reusable.
\section{What This Means: Protocol Libraries}
\label{sec:org5554b17}

Composition through types means session types become \textbf{building blocks}. A
protocol library is a collection of named session types, each describing a
reusable phase of interaction:

\begin{verbatim}
;; --- Protocol building blocks ---

session Handshake
  ! Version . ? Version . Authenticated

session Authenticated
  ! Credentials . ? [Result Token AuthError] .
  +> | :ok -> Ready
     | :fail -> end

;; --- Domain protocols built from blocks ---

session Ready
  +>
    | :query -> QueryMode
    | :close -> end

session SecureDbSession
  Handshake               ;; starts with handshake → authenticated → ready
  ;; (the chain follows from Handshake's continuation declarations)
\end{verbatim}

Each block is independently testable, independently verifiable, and
independently documentable. Complex protocols are assembled from simple,
well-understood components.
\section{The State Machine Interpretation}
\label{sec:org8db9adc}

Named session types are states. Operations within a session type are
transitions. The continuation after each operation names the next state.

\begin{verbatim}
               Handshake
                  │
         ! Version . ? Version
                  │
                  ▼
            Authenticated
                  │
       ! Credentials . ? Result
                  │
             ┌────┴────┐
             ▼         ▼
          Ready      end
           │
    ┌──────┼──────┐
    ▼      ▼      ▼
QueryMode  │    end
    │      │
    └──────┘
 (returns to Ready)
\end{verbatim}

This is a typed, verified finite state machine. Every transition is
type-checked. Dead states are impossible (the type checker rejects protocols
that don't reach \texttt{end}). Invalid transitions are type errors. The session type
\textbf{is} the state machine specification.
\section{The \texttt{end} Question: Intentional Design, Not Missing Feature}
\label{sec:org59a0019}

A natural question arises: if \texttt{Handshake} ends with \texttt{end}, can we compose it
with \texttt{Auth} by replacing \texttt{end} with \texttt{Auth}?

\begin{verbatim}
;; Both are complete, terminal protocols:
session Handshake
  ! Version . ? Version . end

session Auth
  ! Credentials . ? Token . end

;; Can we write this?
;; session Composed = Handshake . Auth    ;; hypothetical end-splicing
\end{verbatim}

This would require \textbf{end-splicing} --- a type-level operation that substitutes
\texttt{end} in one protocol with the start of another. Prologos \textbf{deliberately does
not do this}, for principled reasons.
\subsection{Why End-Splicing Is Rejected}
\label{sec:org81e1389}

End-splicing introduces semantic complications:

\begin{enumerate}
\item \textbf{Multiple exit points}: A protocol with several branches reaching \texttt{end}
would have all of them spliced. Is that always what the author intended?

\item \textbf{Recursive protocol interaction}: If a protocol has recursive
self-references (\texttt{rec}) and \texttt{end} branches, splicing \texttt{end} must create a
new recursive type where the self-references point to the spliced version.
This requires type-level rewriting of recursive bindings.

\item \textbf{Mixed continuations}: If some branches end with \texttt{end} and others with
named types, which \texttt{end}​s get spliced? The semantics become ambiguous.

\item \textbf{Retroactive composition}: A protocol defined with \texttt{end} was designed as
terminal. Retroactively composing it changes its meaning without changing
its definition. This violates referential transparency at the type level.
\end{enumerate}
\subsection{The Intentionality Principle}
\label{sec:orgfcd196b}

When \texttt{Handshake} says \texttt{end}, it means "\emph{I am done --- the connection is
finished}." When \texttt{Handshake} says \texttt{. Authenticated}, it means "\emph{I am a phase,
and after me comes authentication}." These are different protocols with
different meanings.

A handshake-then-close and a handshake-then-authenticate are not the same
handshake with different suffixes. They are different commitments. The
named-continuation approach makes protocol composition \textbf{intentional}:

\begin{itemize}
\item The \textbf{author} of a session type declares what comes next
\item A protocol defined with \texttt{end} is \textbf{complete} --- it is not a composable prefix
\item A protocol defined with a named continuation is \textbf{open} --- it transitions
to another protocol
\item This distinction is explicit, visible in the type definition, and
enforced by the type checker
\end{itemize}

This is analogous to the difference between a function that returns \texttt{A} and a
function that takes a continuation \texttt{(A -> B) -> B}. Both are useful; neither
is the other.
\subsection{Making Protocols Composable}
\label{sec:orgfc697e3}

If a protocol author intends their protocol to be used as a composable phase,
they define it with a named continuation:

\begin{verbatim}
;; Reusable: author explicitly says "something comes after me"
session Handshake
  ! Version . ? Version . Authenticated     ;; open: transitions to Authenticated

;; Terminal: author explicitly says "I'm the end"
session HandshakeDone
  ! Version . ? Version . end               ;; closed: terminates
\end{verbatim}

Two types, two intentions, both clear. No machinery needed to distinguish them.
The author's intent is visible in the type definition.
\section{Capability Composition}
\label{sec:org8a0e235}

When protocols compose, capabilities compose with them. Each named session type
can carry capability requirements. The composed protocol's requirements are the
union of its constituents':

\begin{verbatim}
;; Each phase has its own capability requirements
session Handshake {net :0 TlsCap}
  ! ClientHello . ? ServerHello . Authenticated

session Authenticated {net :0 AuthCap}
  ! Credentials . ? Token . Ready

session Ready {fs :0 ReadCap}
  +>
    | :read -> ! Path . ? String . Ready
    | :done -> end
\end{verbatim}

A process implementing the full \texttt{Handshake → Authenticated → Ready} chain
requires \texttt{\{net :0 TlsCap, net :0 AuthCap, fs :0 ReadCap\}} --- the union of
all phase requirements. The compiler computes this automatically. At runtime,
\texttt{:0} capabilities are erased --- zero cost for authority verification.

This means:
\begin{itemize}
\item \textbf{Protocols compose} --- named types in continuation position
\item \textbf{Capabilities compose} --- union of per-phase requirements
\item \textbf{Types verify the whole thing} --- the compiler checks authority at every
phase transition
\item \textbf{Runtime cost is zero} --- \texttt{:0} authority proofs are erased
\end{itemize}
\section{Relationship to Existing Type System Features}
\label{sec:orgfcf0b2f}

Protocol composition through types is not an isolated feature. It connects
to and reinforces several existing Prologos design principles:
\subsection{Session Types as Layer 5}
\label{sec:org16fa2c1}

The Layered Architecture (DESIGN\textsubscript{PRINCIPLES} §Layered Architecture) places
protocols at Layer 5 (\texttt{session}) and processes at Layer 6 (\texttt{defproc}). Protocol
composition works entirely within Layer 5 --- it is a property of protocol
specifications, independent of process implementation. A \texttt{defproc} implements
one endpoint of a (possibly composed) protocol; it does not need to know
whether the protocol was assembled from multiple named session types.
\subsection{Decomplection of Protocol and Process}
\label{sec:org051ff39}

DESIGN\textsubscript{PRINCIPLES} §Decomplection establishes that protocol specification
(\texttt{session}) is decoupled from process implementation (\texttt{defproc}). Protocol
composition deepens this decomplection: complex protocols can be assembled
from simple building blocks at the type level, without any corresponding
complexity in the process implementation. The process sees a single session
type; whether it was composed from three named types or defined monolithically
is invisible.
\subsection{Duality Composes}
\label{sec:org061ba85}

Session type duality is preserved through composition. If \texttt{Handshake}
transitions to \texttt{Auth}, then \texttt{dual Handshake} transitions to \texttt{dual Auth}. The
type checker verifies duality for the entire composed protocol, not just each
phase independently. This ensures end-to-end protocol safety.
\subsection{Progressive Disclosure}
\label{sec:org11a5e1c}

Protocol composition is a Tier 3 feature --- expert users building complex
concurrent systems. Tier 1 users (\texttt{defn main}) never see session types. Tier 2
users (\texttt{with-open}) interact with single, monolithic protocols. Only users who
need multi-phase, multi-protocol architectures encounter composition. The
complexity is opt-in, following the Disappearing Features principle
(ERGONOMICS §Guiding Principle).
\section{Theoretical Foundations}
\label{sec:org52cf036}

Protocol composition through types has roots in several lines of research,
but Prologos's specific combination is novel.
\subsection{Curry-Howard for Linear Logic (Toninho, Caires, Pfenning, 2011)}
\label{sec:org3fa04f8}

Session types correspond to linear logic propositions under the Curry-Howard
correspondence. Protocol composition corresponds to the \textbf{cut rule} --- the
most fundamental operation in logic. Sequencing protocols is logically
equivalent to composing proofs. The theoretical foundation for what Prologos
does is as old as linear logic itself.
\subsection{Polymorphic Session Types (Caires et al., 2013)}
\label{sec:org3334997}

Session types parameterized by other types, including continuation types. This
gives the formal machinery for open-ended protocols:
\texttt{Handshake[K] = ! Version . ? Version . K}. Prologos achieves the same
effect through named continuations without requiring explicit type parameters.
\subsection{Context-Free Session Types (Thiemann \& Vasconcelos, 2016; Padovani, 2017)}
\label{sec:org496050b}

Extensions that allow session types to express context-free languages, not
just regular languages. This gives more expressive composition at the cost of
decidability for some properties. Prologos stays within the regular fragment
(tail-recursive protocols) for decidability guarantees.
\subsection{Scribble (Honda, Yoshida, et al.)}
\label{sec:orga195220}

A protocol description language for multiparty session types that supports
protocol modules. The closest existing practical tool to protocol composition.
But Scribble is a specification language --- protocols are described, not typed.
Prologos makes protocols first-class types in the programming language itself.
\subsection{What Is Novel}
\label{sec:orgc6e0047}

The \textbf{combination} is novel:

\begin{enumerate}
\item Named session types as composable building blocks within a programming
language (not a separate specification language)
\item Capability composition tracking alongside protocol composition
\item Propagator-based runtime that schedules composed protocols efficiently
\item Progressive disclosure --- composition is opt-in, not forced
\item The intentionality principle --- \texttt{end} means terminal, named continuations
mean composable, and the distinction is a design choice
\end{enumerate}

No existing language or tool unifies all five. The individual pieces are
well-understood; the synthesis is new.
\section{Design Principles for Protocol Composition}
\label{sec:org599fa57}

\begin{enumerate}
\item \textbf{Protocols Are Types.} A session type is a first-class type. It can appear
anywhere a type can appear --- as a continuation, as a type parameter, as a
constraint. Protocol composition is type composition.

\item \textbf{Composition Is Intentional.} The author of a session type decides whether
it is composable (named continuation) or terminal (\texttt{end}). Composition is
not retroactive. The type definition is a commitment.

\item \textbf{Named Continuations, Not End-Splicing.} Protocols compose through named
types in continuation position. There is no mechanism to replace \texttt{end} in a
complete protocol with a continuation. The distinction between open and
closed protocols is part of the protocol's meaning.

\item \textbf{Capabilities Compose With Protocols.} When protocols compose, their
capability requirements compose via union. The compiler tracks authority
through every phase transition. Authority verification is free (\texttt{:0} erased).

\item \textbf{Duality Is Preserved.} Protocol composition preserves the duality
invariant. If server-side composes phases A → B → C, the client-side sees
dual(A) → dual(B) → dual(C). End-to-end safety is guaranteed.

\item \textbf{State Machines Are Types.} The protocol graph (named session types as
states, operations as transitions) is verified by the type checker. Dead
states, invalid transitions, and stuck channels are type errors.

\item \textbf{Progressive Disclosure.} Simple protocols are simple to define. Composition
is available when needed but never required. A monolithic \texttt{session} that
does everything inline is always valid.
\end{enumerate}
\section{Relationship to Other Principles}
\label{sec:org905386c}

\begin{center}
\begin{tabular}{ll}
Document & Relationship\\
\hline
\texttt{DESIGN\_PRINCIPLES.org} & "Correctness Through Types" --- protocol composition is type composition. "Decomplection" --- protocol spec decoupled from process impl. "Layered Architecture" --- composition operates at Layer 5.\\
\texttt{LANGUAGE\_VISION.org} & "Protocol correctness" (§What Problem Does Prologos Solve?) extends from single protocols to composed protocols. "Session Types as Protocol Specifications" (§Cutting-Edge Research) gains composability.\\
\texttt{CAPABILITY\_SECURITY.md} & "Session Types and Capability Protocols" (§Session Types) --- capability delegation follows composed protocol structure. Capability requirements union across composed phases.\\
\texttt{ERGONOMICS.org} & "Disappearing Features" --- composition is Tier 3, invisible to beginners. Progressive disclosure governs when users encounter it.\\
\texttt{LANGUAGE\_DESIGN.org} & "Session Types as Protocol Specifications" (§Session Types) --- composition extends the specification power of session types. Duality, recursion, and choice all compose naturally.\\
\texttt{PATTERNS\_AND\_CONVENTIONS.org} & Session type naming patterns extend to composed protocols. Each named phase follows the same naming conventions.\\
IO Library Design V2 & The IO design uses protocol composition extensively: \texttt{Handshake → Auth → FileOps} patterns, branching into sub-protocols, mode transitions.\\
\end{tabular}
\end{center}
\section{Implications for the Language}
\label{sec:org8fb1dbc}

\subsection{Protocol Libraries Become Possible}
\label{sec:org05ba1fc}

Named, composable session types enable a \textbf{standard library of protocols} ---
reusable communication patterns that users can compose into application-specific
protocols. Examples:

\begin{itemize}
\item \texttt{Handshake} --- version negotiation
\item \texttt{TlsHandshake} --- TLS setup
\item \texttt{Auth} --- credential exchange
\item \texttt{Heartbeat} --- keepalive with timeout
\item \texttt{Pagination} --- cursor-based result traversal
\end{itemize}

These would live in \texttt{prologos.core.protocols} and compose freely.
\subsection{Testing Protocols in Isolation}
\label{sec:org76da910}

Each named session type can be tested independently. A test for \texttt{QueryMode}
doesn't need to set up the full \texttt{DatabaseSession → QueryMode} chain --- it
directly tests the \texttt{QueryMode} protocol against a mock endpoint. This mirrors
how pure functions are tested independently of their callers.
\subsection{Error Handling at Phase Boundaries}
\label{sec:org1480930}

Errors at phase transitions --- authentication failure, handshake mismatch ---
are type-level events. A branch to \texttt{end} at the wrong phase is visible in the
type. Error protocols can be composed just like success protocols:

\begin{verbatim}
session Authenticated
  ! Credentials . ? [Result Token AuthError] .
  +>
    | :ok   -> Ready
    | :fail -> ErrorRecovery           ;; transition to error protocol

session ErrorRecovery
  ! RetryRequest . ? [Result Token AuthError] .
  +>
    | :ok    -> Ready
    | :fail  -> end                     ;; final failure
\end{verbatim}
\subsection{AI Agent Protocols}
\label{sec:org8d8408a}

Multi-agent communication (LANGUAGE\textsubscript{VISION} §AI Agent Infrastructure) benefits
directly from protocol composition. Agent interaction protocols are built from
composable phases:

\begin{verbatim}
session AgentNegotiation
  ! Proposal . ? CounterProposal .
  +>
    | :accept -> CollaborativeWork
    | :counter -> AgentNegotiation     ;; re-negotiate
    | :reject -> end

session CollaborativeWork
  rec
    +>
      | :delegate -> ! Task . ? Result . CollaborativeWork
      | :report   -> ! Summary . end
\end{verbatim}

The session type guarantees that agents follow the negotiation protocol, that
delegation respects authority boundaries (via capabilities), and that the
collaboration terminates.
\section{Inspirations}
\label{sec:orga29a86d}

\begin{itemize}
\item \textbf{Toninho, Caires, Pfenning} --- Session types as propositions (Curry-Howard
for linear logic). Protocol composition = cut rule. The deepest theoretical
foundation for this work.
\item \textbf{Honda, Yoshida, Carbone} --- Multiparty session types and the Scribble
protocol description language. The closest practical precedent for protocol
modularity.
\item \textbf{Lindley \& Morris} --- Semantics for propositions as sessions. The cut rule
gives composition; our named continuations are a user-facing expression
of the same operation.
\item \textbf{Padovani} --- Context-free session types. Shows that session types can
express more than regular languages, though we stay within the regular
fragment for decidability.
\item \textbf{Mark Miller (E language)} --- Object-capability model. Capabilities compose
through delegation; our capability composition through protocol composition
mirrors this at the type level.
\end{itemize}
\end{document}
